\documentclass[11pt]{article}
\usepackage[margin=0.75in]{geometry}
\usepackage{amsmath, amssymb}
\usepackage{booktabs}
\usepackage{enumitem}
\usepackage{graphicx}

\title{\textbf{AP Physics C Lab: Rotational Inertia}}
\author{Aaryan Manrao, Alexander Marynchev, Elijah Park, Joshua Tom, Maximus Neri}
\date{November 19, 2026}

\begin{document}

\maketitle

\section*{Purpose}
The purpose of this lab is to determine how the distribution of the mass of an object affects its moment of inertia. You will be employing a rotational inertia demonstrator with adjustable masses. By exerting a constant torque and adjusting the position of the masses, you will be able to determine the required relationship.

\section*{Materials}
\begin{itemize}[noitemsep]
    \item 1 Rotational Inertia Demonstrator
    \item 4 adjustable masses
    \item 1 hanging mass
    \item String
    \item Ruler
    \item Caliper
    \item Timer
\end{itemize}

\section*{Procedure}
\begin{enumerate}[noitemsep]
    \item Thread the string through one of the holes in the central mandrel and wrap the string around the mandrel, leaving enough string on the end to suspend the mass.
    \item Determine the perpendicular distance from the string to the rotational axis using the caliper.
    \item Attach the mass to the free end of the string and make sure the string is taut.
    \item Measure the distance from the bottom of the mass to the lab table.
    \item Release the system, allowing the mass to drop and the demonstrator to rotate. With the timer, measure the time it takes the mass to reach the table.
    \item Determine the linear acceleration of the mass. From this, determine the tension in the string and the torque on the demonstrator.
    \item Determine the angular acceleration of the demonstrator as the mass dropped. From this, determine the moment of inertia of the demonstrator.
    \item Place the four adjustable masses onto the demonstrator so that they are as close to the rotational axis as possible. Determine the distance from the axis to the center of the four masses.
    \item Repeat steps 3--7. Determine the moment of inertia of the four masses by subtracting the moment of inertia of the demonstrator from your result.
    \item Adjust the adjustable masses so they are now one centimeter further away from the axis of rotation. Repeat the experiment.
    \item Continue to adjust the masses so that they progressively get one centimeter farther from the rotational axis and repeat the experiment each time, until the masses reach the free end of the rods they are attached to.
    \item From your data, generate a graph of moment of inertia (y-axis) and radial distance of the adjustable masses from the axis of rotation (x-axis). Determine the mathematical relationship between the two.
\end{enumerate}

\section*{Raw Data}
\begin{itemize}[noitemsep]
    \item \textbf{Hanging mass ($m$):} $0.2\text{ kg}$ (infinite significant figures as it is a given mass)
    \item \textbf{Mandrel diameter ($D$):} $0.0385 \pm 0.0005\text{ m}$
    \item \textbf{Axis to String ($r = D/2$):} $0.01925 \pm 0.00025\text{ m}$
\end{itemize}

\begin{table}[h!]
\centering
\small
\begin{tabular}{cccc}
\toprule
\textbf{Trial} & \textbf{Adjustable Mass Distance from Axis (m)} & \textbf{Mass Height (m)} & \textbf{Drop Time (s)} \\
\midrule
1 & No Adjustable Mass & $0.590 \pm 0.001$ & $4.47 \pm 0.5$ \\
2 & $0.122 \pm 0.001$ & $0.590 \pm 0.001$ & $5.39 \pm 0.5$ \\
3 & $0.132 \pm 0.001$ & $0.590 \pm 0.001$ & $5.82 \pm 0.5$ \\
4 & $0.142 \pm 0.001$ & $0.590 \pm 0.001$ & $6.26 \pm 0.5$ \\
5 & $0.152 \pm 0.001$ & $0.590 \pm 0.001$ & $6.81 \pm 0.5$ \\
6 & $0.162 \pm 0.001$ & $0.590 \pm 0.001$ & $6.66 \pm 0.5$ \\
7 & $0.172 \pm 0.001$ & $0.590 \pm 0.001$ & $6.93 \pm 0.5$ \\
8 & $0.182 \pm 0.001$ & $0.590 \pm 0.001$ & $7.51 \pm 0.5$ \\
9 & $0.192 \pm 0.001$ & $0.590 \pm 0.001$ & $7.89 \pm 0.5$ \\
10 & $0.202 \pm 0.001$ & $0.590 \pm 0.001$ & $8.12 \pm 0.5$ \\
11 & $0.212 \pm 0.001$ & $0.590 \pm 0.001$ & $8.37 \pm 0.5$ \\
12 & $0.222 \pm 0.001$ & $0.590 \pm 0.001$ & $8.83 \pm 0.5$ \\
13 & $0.232 \pm 0.001$ & $0.590 \pm 0.001$ & $9.06 \pm 0.5$ \\
14 & $0.242 \pm 0.001$ & $0.590 \pm 0.001$ & $9.36 \pm 0.5$ \\
15 & $0.252 \pm 0.001$ & $0.590 \pm 0.001$ & $9.47 \pm 0.5$ \\
16 & $0.262 \pm 0.001$ & $0.590 \pm 0.001$ & $9.91 \pm 0.5$ \\
17 & $0.272 \pm 0.001$ & $0.590 \pm 0.001$ & $10.40 \pm 0.5$ \\
18 & $0.282 \pm 0.001$ & $0.590 \pm 0.001$ & $10.76 \pm 0.5$ \\
19 & $0.292 \pm 0.001$ & $0.590 \pm 0.001$ & $11.12 \pm 0.5$ \\
20 & $0.302 \pm 0.001$ & $0.590 \pm 0.001$ & $12.24 \pm 0.5$ \\
21 & $0.312 \pm 0.001$ & $0.590 \pm 0.001$ & $13.64 \pm 0.5$ \\
\bottomrule
\end{tabular}
\end{table}

\section*{Analysis}

\textbf{Trial 1 — Demonstrator Only (No Adjustable Masses):}
\begin{align*}
a &= \frac{2(0.590)}{(4.47)^2} = 0.0591\text{ m/s}^2 \\
\delta a &= 0.0591\sqrt{\left(\frac{0.001}{0.590}\right)^2 + \left(\frac{2(0.5)}{4.47}\right)^2} = 0.0132\text{ m/s}^2 \\
g - a &= 9.807 - 0.0591 = 9.748 \pm 0.013\text{ m/s}^2 \\
T &= 0.2(9.748) = 1.950\text{ N}, \qquad \delta T = 0.2(0.0132) = 0.003\text{ N} \\
\tau &= (1.950)(0.01925) = 0.0375\text{ N}\cdot\text{m} \\
\delta\tau &= 0.0375\sqrt{\left(\frac{0.003}{1.950}\right)^2 + \left(\frac{0.00025}{0.01925}\right)^2} = 0.0005\text{ N}\cdot\text{m} \\
\alpha &= \frac{0.0591}{0.01925} = 3.07\text{ rad/s}^2 \\
\delta\alpha &= 3.07\sqrt{\left(\frac{0.0132}{0.0591}\right)^2 + \left(\frac{0.00025}{0.01925}\right)^2} = 0.69\text{ rad/s}^2 \\
I &= \frac{0.0375}{3.07} = 0.0122\text{ kg}\cdot\text{m}^2 \\
\delta I &= 0.0122\sqrt{\left(\frac{0.0005}{0.0375}\right)^2 + \left(\frac{0.69}{3.07}\right)^2} = 0.0027\text{ kg}\cdot\text{m}^2 \\
I_{\text{demo}} &= 0.0122 \pm 0.0027\text{ kg}\cdot\text{m}^2
\end{align*}

\textbf{Trial 2 — $R = 0.122 \pm 0.001\text{ m}$:}
\begin{align*}
a &= \frac{2(0.590)}{(5.39)^2} = 0.041\text{ m/s}^2 \\
\delta a &= 0.041\sqrt{\left(\frac{0.001}{0.590}\right)^2 + \left(\frac{2(0.5)}{5.39}\right)^2} = 0.008\text{ m/s}^2 \\
g - a &= 9.807 - 0.041 = 9.766 \pm 0.008\text{ m/s}^2 \\
T &= 0.2(9.766) = 1.953\text{ N}, \qquad \delta T = 0.2(0.008) = 0.002\text{ N} \\
\tau &= (1.953)(0.01925) = 0.0376\text{ N}\cdot\text{m} \\
\delta\tau &= 0.0376\sqrt{\left(\frac{0.002}{1.953}\right)^2 + \left(\frac{0.00025}{0.01925}\right)^2} = 0.0005\text{ N}\cdot\text{m} \\
\alpha &= \frac{0.041}{0.01925} = 2.11\text{ rad/s}^2 \\
\delta\alpha &= 2.11\sqrt{\left(\frac{0.008}{0.041}\right)^2 + \left(\frac{0.00025}{0.01925}\right)^2} = 0.39\text{ rad/s}^2 \\
I &= \frac{0.0376}{2.11} = 0.018\text{ kg}\cdot\text{m}^2 \\
\delta I &= 0.018\sqrt{\left(\frac{0.0005}{0.0376}\right)^2 + \left(\frac{0.39}{2.11}\right)^2} = 0.003\text{ kg}\cdot\text{m}^2 \\
I_{\text{masses}} &= 0.018 - 0.0122 = 0.006\text{ kg}\cdot\text{m}^2 \\
\delta I_{\text{masses}} &= \sqrt{0.003^2 + 0.0027^2} = 0.004\text{ kg}\cdot\text{m}^2
\end{align*}

\textbf{Trial 3 — $R = 0.132 \pm 0.001\text{ m}$:}
\begin{align*}
a &= \frac{2(0.590)}{(5.82)^2} = 0.035\text{ m/s}^2 \\
\delta a &= 0.035\sqrt{\left(\frac{0.001}{0.590}\right)^2 + \left(\frac{2(0.5)}{5.82}\right)^2} = 0.006\text{ m/s}^2 \\
g - a &= 9.807 - 0.035 = 9.772 \pm 0.006\text{ m/s}^2 \\
T &= 0.2(9.772) = 1.954\text{ N}, \qquad \delta T = 0.2(0.006) = 0.001\text{ N} \\
\tau &= (1.954)(0.01925) = 0.0376\text{ N}\cdot\text{m} \\
\delta\tau &= 0.0376\sqrt{\left(\frac{0.001}{1.954}\right)^2 + \left(\frac{0.00025}{0.01925}\right)^2} = 0.0005\text{ N}\cdot\text{m} \\
\alpha &= \frac{0.035}{0.01925} = 1.81\text{ rad/s}^2 \\
\delta\alpha &= 1.81\sqrt{\left(\frac{0.006}{0.035}\right)^2 + \left(\frac{0.00025}{0.01925}\right)^2} = 0.31\text{ rad/s}^2 \\
I &= \frac{0.0376}{1.81} = 0.021\text{ kg}\cdot\text{m}^2 \\
\delta I &= 0.021\sqrt{\left(\frac{0.0005}{0.0376}\right)^2 + \left(\frac{0.31}{1.81}\right)^2} = 0.004\text{ kg}\cdot\text{m}^2 \\
I_{\text{masses}} &= 0.021 - 0.0122 = 0.009\text{ kg}\cdot\text{m}^2 \\
\delta I_{\text{masses}} &= \sqrt{0.004^2 + 0.0027^2} = 0.005\text{ kg}\cdot\text{m}^2
\end{align*}

\textbf{Trial 4 — $R = 0.142 \pm 0.001\text{ m}$:}
\begin{align*}
a &= \frac{2(0.590)}{(6.26)^2} = 0.030\text{ m/s}^2 \\
\delta a &= 0.030\sqrt{\left(\frac{0.001}{0.590}\right)^2 + \left(\frac{2(0.5)}{6.26}\right)^2} = 0.005\text{ m/s}^2 \\
g - a &= 9.807 - 0.030 = 9.777 \pm 0.005\text{ m/s}^2 \\
T &= 0.2(9.777) = 1.955\text{ N}, \qquad \delta T = 0.2(0.005) = 0.001\text{ N} \\
\tau &= (1.955)(0.01925) = 0.0376\text{ N}\cdot\text{m} \\
\delta\tau &= 0.0376\sqrt{\left(\frac{0.001}{1.955}\right)^2 + \left(\frac{0.00025}{0.01925}\right)^2} = 0.0005\text{ N}\cdot\text{m} \\
\alpha &= \frac{0.030}{0.01925} = 1.56\text{ rad/s}^2 \\
\delta\alpha &= 1.56\sqrt{\left(\frac{0.005}{0.030}\right)^2 + \left(\frac{0.00025}{0.01925}\right)^2} = 0.25\text{ rad/s}^2 \\
I &= \frac{0.0376}{1.56} = 0.024\text{ kg}\cdot\text{m}^2 \\
\delta I &= 0.024\sqrt{\left(\frac{0.0005}{0.0376}\right)^2 + \left(\frac{0.25}{1.56}\right)^2} = 0.004\text{ kg}\cdot\text{m}^2 \\
I_{\text{masses}} &= 0.024 - 0.0122 = 0.012\text{ kg}\cdot\text{m}^2 \\
\delta I_{\text{masses}} &= \sqrt{0.004^2 + 0.0027^2} = 0.005\text{ kg}\cdot\text{m}^2
\end{align*}

\textbf{Trial 5 — $R = 0.152 \pm 0.001\text{ m}$:}
\begin{align*}
a &= \frac{2(0.590)}{(6.81)^2} = 0.025\text{ m/s}^2 \\
\delta a &= 0.025\sqrt{\left(\frac{0.001}{0.590}\right)^2 + \left(\frac{2(0.5)}{6.81}\right)^2} = 0.004\text{ m/s}^2 \\
g - a &= 9.807 - 0.025 = 9.782 \pm 0.004\text{ m/s}^2 \\
T &= 0.2(9.782) = 1.956\text{ N}, \qquad \delta T = 0.2(0.004) = 0.001\text{ N} \\
\tau &= (1.956)(0.01925) = 0.0377\text{ N}\cdot\text{m} \\
\delta\tau &= 0.0377\sqrt{\left(\frac{0.001}{1.956}\right)^2 + \left(\frac{0.00025}{0.01925}\right)^2} = 0.0005\text{ N}\cdot\text{m} \\
\alpha &= \frac{0.025}{0.01925} = 1.32\text{ rad/s}^2 \\
\delta\alpha &= 1.32\sqrt{\left(\frac{0.004}{0.025}\right)^2 + \left(\frac{0.00025}{0.01925}\right)^2} = 0.19\text{ rad/s}^2 \\
I &= \frac{0.0377}{1.32} = 0.028\text{ kg}\cdot\text{m}^2 \\
\delta I &= 0.028\sqrt{\left(\frac{0.0005}{0.0377}\right)^2 + \left(\frac{0.19}{1.32}\right)^2} = 0.004\text{ kg}\cdot\text{m}^2 \\
I_{\text{masses}} &= 0.028 - 0.0122 = 0.016\text{ kg}\cdot\text{m}^2 \\
\delta I_{\text{masses}} &= \sqrt{0.004^2 + 0.0027^2} = 0.005\text{ kg}\cdot\text{m}^2
\end{align*}

\textbf{Trial 6 — $R = 0.162 \pm 0.001\text{ m}$:}
\begin{align*}
a &= \frac{2(0.590)}{(6.66)^2} = 0.027\text{ m/s}^2 \\
\delta a &= 0.027\sqrt{\left(\frac{0.001}{0.590}\right)^2 + \left(\frac{2(0.5)}{6.66}\right)^2} = 0.004\text{ m/s}^2 \\
g - a &= 9.807 - 0.027 = 9.780 \pm 0.004\text{ m/s}^2 \\
T &= 0.2(9.780) = 1.956\text{ N}, \qquad \delta T = 0.2(0.004) = 0.001\text{ N} \\
\tau &= (1.956)(0.01925) = 0.0377\text{ N}\cdot\text{m} \\
\delta\tau &= 0.0377\sqrt{\left(\frac{0.001}{1.956}\right)^2 + \left(\frac{0.00025}{0.01925}\right)^2} = 0.0005\text{ N}\cdot\text{m} \\
\alpha &= \frac{0.027}{0.01925} = 1.38\text{ rad/s}^2 \\
\delta\alpha &= 1.38\sqrt{\left(\frac{0.004}{0.027}\right)^2 + \left(\frac{0.00025}{0.01925}\right)^2} = 0.21\text{ rad/s}^2 \\
I &= \frac{0.0377}{1.38} = 0.027\text{ kg}\cdot\text{m}^2 \\
\delta I &= 0.027\sqrt{\left(\frac{0.0005}{0.0377}\right)^2 + \left(\frac{0.21}{1.38}\right)^2} = 0.004\text{ kg}\cdot\text{m}^2 \\
I_{\text{masses}} &= 0.027 - 0.0122 = 0.015\text{ kg}\cdot\text{m}^2 \\
\delta I_{\text{masses}} &= \sqrt{0.004^2 + 0.0027^2} = 0.005\text{ kg}\cdot\text{m}^2
\end{align*}

\textbf{Trial 7 — $R = 0.172 \pm 0.001\text{ m}$:}
\begin{align*}
a &= \frac{2(0.590)}{(6.93)^2} = 0.025\text{ m/s}^2 \\
\delta a &= 0.025\sqrt{\left(\frac{0.001}{0.590}\right)^2 + \left(\frac{2(0.5)}{6.93}\right)^2} = 0.004\text{ m/s}^2 \\
g - a &= 9.807 - 0.025 = 9.782 \pm 0.004\text{ m/s}^2 \\
T &= 0.2(9.782) = 1.956\text{ N}, \qquad \delta T = 0.2(0.004) = 0.001\text{ N} \\
\tau &= (1.956)(0.01925) = 0.0377\text{ N}\cdot\text{m} \\
\delta\tau &= 0.0377\sqrt{\left(\frac{0.001}{1.956}\right)^2 + \left(\frac{0.00025}{0.01925}\right)^2} = 0.0005\text{ N}\cdot\text{m} \\
\alpha &= \frac{0.025}{0.01925} = 1.28\text{ rad/s}^2 \\
\delta\alpha &= 1.28\sqrt{\left(\frac{0.004}{0.025}\right)^2 + \left(\frac{0.00025}{0.01925}\right)^2} = 0.18\text{ rad/s}^2 \\
I &= \frac{0.0377}{1.28} = 0.030\text{ kg}\cdot\text{m}^2 \\
\delta I &= 0.030\sqrt{\left(\frac{0.0005}{0.0377}\right)^2 + \left(\frac{0.18}{1.28}\right)^2} = 0.004\text{ kg}\cdot\text{m}^2 \\
I_{\text{masses}} &= 0.030 - 0.0122 = 0.018\text{ kg}\cdot\text{m}^2 \\
\delta I_{\text{masses}} &= \sqrt{0.004^2 + 0.0027^2} = 0.005\text{ kg}\cdot\text{m}^2
\end{align*}

\textbf{Trial 8 — $R = 0.182 \pm 0.001\text{ m}$:}
\begin{align*}
a &= \frac{2(0.590)}{(7.51)^2} = 0.021\text{ m/s}^2 \\
\delta a &= 0.021\sqrt{\left(\frac{0.001}{0.590}\right)^2 + \left(\frac{2(0.5)}{7.51}\right)^2} = 0.003\text{ m/s}^2 \\
g - a &= 9.807 - 0.021 = 9.786 \pm 0.003\text{ m/s}^2 \\
T &= 0.2(9.786) = 1.957\text{ N}, \qquad \delta T = 0.2(0.003) = 0.001\text{ N} \\
\tau &= (1.957)(0.01925) = 0.0377\text{ N}\cdot\text{m} \\
\delta\tau &= 0.0377\sqrt{\left(\frac{0.001}{1.957}\right)^2 + \left(\frac{0.00025}{0.01925}\right)^2} = 0.0005\text{ N}\cdot\text{m} \\
\alpha &= \frac{0.021}{0.01925} = 1.09\text{ rad/s}^2 \\
\delta\alpha &= 1.09\sqrt{\left(\frac{0.003}{0.021}\right)^2 + \left(\frac{0.00025}{0.01925}\right)^2} = 0.15\text{ rad/s}^2 \\
I &= \frac{0.0377}{1.09} = 0.035\text{ kg}\cdot\text{m}^2 \\
\delta I &= 0.035\sqrt{\left(\frac{0.0005}{0.0377}\right)^2 + \left(\frac{0.15}{1.09}\right)^2} = 0.005\text{ kg}\cdot\text{m}^2 \\
I_{\text{masses}} &= 0.035 - 0.0122 = 0.023\text{ kg}\cdot\text{m}^2 \\
\delta I_{\text{masses}} &= \sqrt{0.005^2 + 0.0027^2} = 0.006\text{ kg}\cdot\text{m}^2
\end{align*}

\textbf{Trial 9 — $R = 0.192 \pm 0.001\text{ m}$:}
\begin{align*}
a &= \frac{2(0.590)}{(7.89)^2} = 0.019\text{ m/s}^2 \\
\delta a &= 0.019\sqrt{\left(\frac{0.001}{0.590}\right)^2 + \left(\frac{2(0.5)}{7.89}\right)^2} = 0.002\text{ m/s}^2 \\
g - a &= 9.807 - 0.019 = 9.788 \pm 0.002\text{ m/s}^2 \\
T &= 0.2(9.788) = 1.958\text{ N}, \qquad \delta T = 0.2(0.002) = 0.000\text{ N} \\
\tau &= (1.958)(0.01925) = 0.0377\text{ N}\cdot\text{m} \\
\delta\tau &= 0.0377\sqrt{\left(\frac{0.000}{1.958}\right)^2 + \left(\frac{0.00025}{0.01925}\right)^2} = 0.0005\text{ N}\cdot\text{m} \\
\alpha &= \frac{0.019}{0.01925} = 0.98\text{ rad/s}^2 \\
\delta\alpha &= 0.98\sqrt{\left(\frac{0.002}{0.019}\right)^2 + \left(\frac{0.00025}{0.01925}\right)^2} = 0.13\text{ rad/s}^2 \\
I &= \frac{0.0377}{0.98} = 0.038\text{ kg}\cdot\text{m}^2 \\
\delta I &= 0.038\sqrt{\left(\frac{0.0005}{0.0377}\right)^2 + \left(\frac{0.13}{0.98}\right)^2} = 0.005\text{ kg}\cdot\text{m}^2 \\
I_{\text{masses}} &= 0.038 - 0.0122 = 0.026\text{ kg}\cdot\text{m}^2 \\
\delta I_{\text{masses}} &= \sqrt{0.005^2 + 0.0027^2} = 0.006\text{ kg}\cdot\text{m}^2
\end{align*}

\textbf{Trial 10 — $R = 0.202 \pm 0.001\text{ m}$:}
\begin{align*}
a &= \frac{2(0.590)}{(8.12)^2} = 0.018\text{ m/s}^2 \\
\delta a &= 0.018\sqrt{\left(\frac{0.001}{0.590}\right)^2 + \left(\frac{2(0.5)}{8.12}\right)^2} = 0.002\text{ m/s}^2 \\
g - a &= 9.807 - 0.018 = 9.789 \pm 0.002\text{ m/s}^2 \\
T &= 0.2(9.789) = 1.958\text{ N}, \qquad \delta T = 0.2(0.002) = 0.000\text{ N} \\
\tau &= (1.958)(0.01925) = 0.0377\text{ N}\cdot\text{m} \\
\delta\tau &= 0.0377\sqrt{\left(\frac{0.000}{1.958}\right)^2 + \left(\frac{0.00025}{0.01925}\right)^2} = 0.0005\text{ N}\cdot\text{m} \\
\alpha &= \frac{0.018}{0.01925} = 0.93\text{ rad/s}^2 \\
\delta\alpha &= 0.93\sqrt{\left(\frac{0.002}{0.018}\right)^2 + \left(\frac{0.00025}{0.01925}\right)^2} = 0.12\text{ rad/s}^2 \\
I &= \frac{0.0377}{0.93} = 0.041\text{ kg}\cdot\text{m}^2 \\
\delta I &= 0.041\sqrt{\left(\frac{0.0005}{0.0377}\right)^2 + \left(\frac{0.12}{0.93}\right)^2} = 0.005\text{ kg}\cdot\text{m}^2 \\
I_{\text{masses}} &= 0.041 - 0.0122 = 0.029\text{ kg}\cdot\text{m}^2 \\
\delta I_{\text{masses}} &= \sqrt{0.005^2 + 0.0027^2} = 0.006\text{ kg}\cdot\text{m}^2
\end{align*}

\textbf{Trial 11 — $R = 0.212 \pm 0.001\text{ m}$:}
\begin{align*}
a &= \frac{2(0.590)}{(8.37)^2} = 0.017\text{ m/s}^2 \\
\delta a &= 0.017\sqrt{\left(\frac{0.001}{0.590}\right)^2 + \left(\frac{2(0.5)}{8.37}\right)^2} = 0.002\text{ m/s}^2 \\
g - a &= 9.807 - 0.017 = 9.790 \pm 0.002\text{ m/s}^2 \\
T &= 0.2(9.790) = 1.958\text{ N}, \qquad \delta T = 0.2(0.002) = 0.000\text{ N} \\
\tau &= (1.958)(0.01925) = 0.0377\text{ N}\cdot\text{m} \\
\delta\tau &= 0.0377\sqrt{\left(\frac{0.000}{1.958}\right)^2 + \left(\frac{0.00025}{0.01925}\right)^2} = 0.0005\text{ N}\cdot\text{m} \\
\alpha &= \frac{0.017}{0.01925} = 0.87\text{ rad/s}^2 \\
\delta\alpha &= 0.87\sqrt{\left(\frac{0.002}{0.017}\right)^2 + \left(\frac{0.00025}{0.01925}\right)^2} = 0.11\text{ rad/s}^2 \\
I &= \frac{0.0377}{0.87} = 0.043\text{ kg}\cdot\text{m}^2 \\
\delta I &= 0.043\sqrt{\left(\frac{0.0005}{0.0377}\right)^2 + \left(\frac{0.11}{0.87}\right)^2} = 0.005\text{ kg}\cdot\text{m}^2 \\
I_{\text{masses}} &= 0.043 - 0.0122 = 0.031\text{ kg}\cdot\text{m}^2 \\
\delta I_{\text{masses}} &= \sqrt{0.005^2 + 0.0027^2} = 0.006\text{ kg}\cdot\text{m}^2
\end{align*}

\textbf{Trial 12 — $R = 0.222 \pm 0.001\text{ m}$:}
\begin{align*}
a &= \frac{2(0.590)}{(8.83)^2} = 0.015\text{ m/s}^2 \\
\delta a &= 0.015\sqrt{\left(\frac{0.001}{0.590}\right)^2 + \left(\frac{2(0.5)}{8.83}\right)^2} = 0.002\text{ m/s}^2 \\
g - a &= 9.807 - 0.015 = 9.792 \pm 0.002\text{ m/s}^2 \\
T &= 0.2(9.792) = 1.958\text{ N}, \qquad \delta T = 0.2(0.002) = 0.000\text{ N} \\
\tau &= (1.958)(0.01925) = 0.0377\text{ N}\cdot\text{m} \\
\delta\tau &= 0.0377\sqrt{\left(\frac{0.000}{1.958}\right)^2 + \left(\frac{0.00025}{0.01925}\right)^2} = 0.0005\text{ N}\cdot\text{m} \\
\alpha &= \frac{0.015}{0.01925} = 0.79\text{ rad/s}^2 \\
\delta\alpha &= 0.79\sqrt{\left(\frac{0.002}{0.015}\right)^2 + \left(\frac{0.00025}{0.01925}\right)^2} = 0.09\text{ rad/s}^2 \\
I &= \frac{0.0377}{0.79} = 0.048\text{ kg}\cdot\text{m}^2 \\
\delta I &= 0.048\sqrt{\left(\frac{0.0005}{0.0377}\right)^2 + \left(\frac{0.09}{0.79}\right)^2} = 0.006\text{ kg}\cdot\text{m}^2 \\
I_{\text{masses}} &= 0.048 - 0.0122 = 0.036\text{ kg}\cdot\text{m}^2 \\
\delta I_{\text{masses}} &= \sqrt{0.006^2 + 0.0027^2} = 0.007\text{ kg}\cdot\text{m}^2
\end{align*}

\textbf{Trial 13 — $R = 0.232 \pm 0.001\text{ m}$:}
\begin{align*}
a &= \frac{2(0.590)}{(9.06)^2} = 0.014\text{ m/s}^2 \\
\delta a &= 0.014\sqrt{\left(\frac{0.001}{0.590}\right)^2 + \left(\frac{2(0.5)}{9.06}\right)^2} = 0.002\text{ m/s}^2 \\
g - a &= 9.807 - 0.014 = 9.793 \pm 0.002\text{ m/s}^2 \\
T &= 0.2(9.793) = 1.959\text{ N}, \qquad \delta T = 0.2(0.002) = 0.000\text{ N} \\
\tau &= (1.959)(0.01925) = 0.0377\text{ N}\cdot\text{m} \\
\delta\tau &= 0.0377\sqrt{\left(\frac{0.000}{1.959}\right)^2 + \left(\frac{0.00025}{0.01925}\right)^2} = 0.0005\text{ N}\cdot\text{m} \\
\alpha &= \frac{0.014}{0.01925} = 0.75\text{ rad/s}^2 \\
\delta\alpha &= 0.75\sqrt{\left(\frac{0.002}{0.014}\right)^2 + \left(\frac{0.00025}{0.01925}\right)^2} = 0.08\text{ rad/s}^2 \\
I &= \frac{0.0377}{0.75} = 0.050\text{ kg}\cdot\text{m}^2 \\
\delta I &= 0.050\sqrt{\left(\frac{0.0005}{0.0377}\right)^2 + \left(\frac{0.08}{0.75}\right)^2} = 0.006\text{ kg}\cdot\text{m}^2 \\
I_{\text{masses}} &= 0.050 - 0.0122 = 0.038\text{ kg}\cdot\text{m}^2 \\
\delta I_{\text{masses}} &= \sqrt{0.006^2 + 0.0027^2} = 0.007\text{ kg}\cdot\text{m}^2
\end{align*}

\textbf{Trial 14 — $R = 0.242 \pm 0.001\text{ m}$:}
\begin{align*}
a &= \frac{2(0.590)}{(9.36)^2} = 0.013\text{ m/s}^2 \\
\delta a &= 0.013\sqrt{\left(\frac{0.001}{0.590}\right)^2 + \left(\frac{2(0.5)}{9.36}\right)^2} = 0.001\text{ m/s}^2 \\
g - a &= 9.807 - 0.013 = 9.794 \pm 0.001\text{ m/s}^2 \\
T &= 0.2(9.794) = 1.959\text{ N}, \qquad \delta T = 0.2(0.001) = 0.000\text{ N} \\
\tau &= (1.959)(0.01925) = 0.0377\text{ N}\cdot\text{m} \\
\delta\tau &= 0.0377\sqrt{\left(\frac{0.000}{1.959}\right)^2 + \left(\frac{0.00025}{0.01925}\right)^2} = 0.0005\text{ N}\cdot\text{m} \\
\alpha &= \frac{0.013}{0.01925} = 0.70\text{ rad/s}^2 \\
\delta\alpha &= 0.70\sqrt{\left(\frac{0.001}{0.013}\right)^2 + \left(\frac{0.00025}{0.01925}\right)^2} = 0.08\text{ rad/s}^2 \\
I &= \frac{0.0377}{0.70} = 0.054\text{ kg}\cdot\text{m}^2 \\
\delta I &= 0.054\sqrt{\left(\frac{0.0005}{0.0377}\right)^2 + \left(\frac{0.08}{0.70}\right)^2} = 0.006\text{ kg}\cdot\text{m}^2 \\
I_{\text{masses}} &= 0.054 - 0.0122 = 0.042\text{ kg}\cdot\text{m}^2 \\
\delta I_{\text{masses}} &= \sqrt{0.006^2 + 0.0027^2} = 0.007\text{ kg}\cdot\text{m}^2
\end{align*}

\textbf{Trial 15 — $R = 0.252 \pm 0.001\text{ m}$:}
\begin{align*}
a &= \frac{2(0.590)}{(9.47)^2} = 0.013\text{ m/s}^2 \\
\delta a &= 0.013\sqrt{\left(\frac{0.001}{0.590}\right)^2 + \left(\frac{2(0.5)}{9.47}\right)^2} = 0.001\text{ m/s}^2 \\
g - a &= 9.807 - 0.013 = 9.794 \pm 0.001\text{ m/s}^2 \\
T &= 0.2(9.794) = 1.959\text{ N}, \qquad \delta T = 0.2(0.001) = 0.000\text{ N} \\
\tau &= (1.959)(0.01925) = 0.0377\text{ N}\cdot\text{m} \\
\delta\tau &= 0.0377\sqrt{\left(\frac{0.000}{1.959}\right)^2 + \left(\frac{0.00025}{0.01925}\right)^2} = 0.0005\text{ N}\cdot\text{m} \\
\alpha &= \frac{0.013}{0.01925} = 0.68\text{ rad/s}^2 \\
\delta\alpha &= 0.68\sqrt{\left(\frac{0.001}{0.013}\right)^2 + \left(\frac{0.00025}{0.01925}\right)^2} = 0.07\text{ rad/s}^2 \\
I &= \frac{0.0377}{0.68} = 0.055\text{ kg}\cdot\text{m}^2 \\
\delta I &= 0.055\sqrt{\left(\frac{0.0005}{0.0377}\right)^2 + \left(\frac{0.07}{0.68}\right)^2} = 0.006\text{ kg}\cdot\text{m}^2 \\
I_{\text{masses}} &= 0.055 - 0.0122 = 0.043\text{ kg}\cdot\text{m}^2 \\
\delta I_{\text{masses}} &= \sqrt{0.006^2 + 0.0027^2} = 0.007\text{ kg}\cdot\text{m}^2
\end{align*}

\textbf{Trial 16 — $R = 0.262 \pm 0.001\text{ m}$:}
\begin{align*}
a &= \frac{2(0.590)}{(9.91)^2} = 0.012\text{ m/s}^2 \\
\delta a &= 0.012\sqrt{\left(\frac{0.001}{0.590}\right)^2 + \left(\frac{2(0.5)}{9.91}\right)^2} = 0.001\text{ m/s}^2 \\
g - a &= 9.807 - 0.012 = 9.795 \pm 0.001\text{ m/s}^2 \\
T &= 0.2(9.795) = 1.959\text{ N}, \qquad \delta T = 0.2(0.001) = 0.000\text{ N} \\
\tau &= (1.959)(0.01925) = 0.0377\text{ N}\cdot\text{m} \\
\delta\tau &= 0.0377\sqrt{\left(\frac{0.000}{1.959}\right)^2 + \left(\frac{0.00025}{0.01925}\right)^2} = 0.0005\text{ N}\cdot\text{m} \\
\alpha &= \frac{0.012}{0.01925} = 0.62\text{ rad/s}^2 \\
\delta\alpha &= 0.62\sqrt{\left(\frac{0.001}{0.012}\right)^2 + \left(\frac{0.00025}{0.01925}\right)^2} = 0.06\text{ rad/s}^2 \\
I &= \frac{0.0377}{0.62} = 0.060\text{ kg}\cdot\text{m}^2 \\
\delta I &= 0.060\sqrt{\left(\frac{0.0005}{0.0377}\right)^2 + \left(\frac{0.06}{0.62}\right)^2} = 0.006\text{ kg}\cdot\text{m}^2 \\
I_{\text{masses}} &= 0.060 - 0.0122 = 0.048\text{ kg}\cdot\text{m}^2 \\
\delta I_{\text{masses}} &= \sqrt{0.006^2 + 0.0027^2} = 0.007\text{ kg}\cdot\text{m}^2
\end{align*}

\textbf{Trial 17 — $R = 0.272 \pm 0.001\text{ m}$:}
\begin{align*}
a &= \frac{2(0.590)}{(10.40)^2} = 0.01091\text{ m/s}^2 \\
\delta a &= 0.01091\sqrt{\left(\frac{0.001}{0.590}\right)^2 + \left(\frac{2(0.5)}{10.40}\right)^2} = 0.00105\text{ m/s}^2 \\
g - a &= 9.807 - 0.01091 = 9.796 \pm 0.001\text{ m/s}^2 \\
T &= 0.2(9.796) = 1.959\text{ N}, \qquad \delta T = 0.2(0.00105) = 0.000\text{ N} \\
\tau &= (1.959)(0.01925) = 0.0377\text{ N}\cdot\text{m} \\
\delta\tau &= 0.0377\sqrt{\left(\frac{0.000}{1.959}\right)^2 + \left(\frac{0.00025}{0.01925}\right)^2} = 0.0005\text{ N}\cdot\text{m} \\
\alpha &= \frac{0.01091}{0.01925} = 0.567\text{ rad/s}^2 \\
\delta\alpha &= 0.567\sqrt{\left(\frac{0.00105}{0.01091}\right)^2 + \left(\frac{0.00025}{0.01925}\right)^2} = 0.055\text{ rad/s}^2 \\
I &= \frac{0.0377}{0.567} = 0.0665\text{ kg}\cdot\text{m}^2 \\
\delta I &= 0.0665\sqrt{\left(\frac{0.0005}{0.0377}\right)^2 + \left(\frac{0.055}{0.567}\right)^2} = 0.0065\text{ kg}\cdot\text{m}^2 \\
I_{\text{masses}} &= 0.0665 - 0.0122 = 0.0543\text{ kg}\cdot\text{m}^2 \\
\delta I_{\text{masses}} &= \sqrt{0.0065^2 + 0.0027^2} = 0.0070\text{ kg}\cdot\text{m}^2
\end{align*}

\textbf{Trial 18 — $R = 0.282 \pm 0.001\text{ m}$:}
\begin{align*}
a &= \frac{2(0.590)}{(10.76)^2} = 0.01019\text{ m/s}^2 \\
\delta a &= 0.01019\sqrt{\left(\frac{0.001}{0.590}\right)^2 + \left(\frac{2(0.5)}{10.76}\right)^2} = 0.00095\text{ m/s}^2 \\
g - a &= 9.807 - 0.01019 = 9.797 \pm 0.001\text{ m/s}^2 \\
T &= 0.2(9.797) = 1.959\text{ N}, \qquad \delta T = 0.2(0.00095) = 0.000\text{ N} \\
\tau &= (1.959)(0.01925) = 0.0377\text{ N}\cdot\text{m} \\
\delta\tau &= 0.0377\sqrt{\left(\frac{0.000}{1.959}\right)^2 + \left(\frac{0.00025}{0.01925}\right)^2} = 0.0005\text{ N}\cdot\text{m} \\
\alpha &= \frac{0.01019}{0.01925} = 0.529\text{ rad/s}^2 \\
\delta\alpha &= 0.529\sqrt{\left(\frac{0.00095}{0.01019}\right)^2 + \left(\frac{0.00025}{0.01925}\right)^2} = 0.050\text{ rad/s}^2 \\
I &= \frac{0.0377}{0.529} = 0.0712\text{ kg}\cdot\text{m}^2 \\
\delta I &= 0.0712\sqrt{\left(\frac{0.0005}{0.0377}\right)^2 + \left(\frac{0.050}{0.529}\right)^2} = 0.0067\text{ kg}\cdot\text{m}^2 \\
I_{\text{masses}} &= 0.0712 - 0.0122 = 0.0590\text{ kg}\cdot\text{m}^2 \\
\delta I_{\text{masses}} &= \sqrt{0.0067^2 + 0.0027^2} = 0.0072\text{ kg}\cdot\text{m}^2
\end{align*}

\textbf{Trial 19 — $R = 0.292 \pm 0.001\text{ m}$:}
\begin{align*}
a &= \frac{2(0.590)}{(11.12)^2} = 0.00954\text{ m/s}^2 \\
\delta a &= 0.00954\sqrt{\left(\frac{0.001}{0.590}\right)^2 + \left(\frac{2(0.5)}{11.12}\right)^2} = 0.00086\text{ m/s}^2 \\
g - a &= 9.807 - 0.00954 = 9.797 \pm 0.001\text{ m/s}^2 \\
T &= 0.2(9.797) = 1.959\text{ N}, \qquad \delta T = 0.2(0.00086) = 0.000\text{ N} \\
\tau &= (1.959)(0.01925) = 0.0377\text{ N}\cdot\text{m} \\
\delta\tau &= 0.0377\sqrt{\left(\frac{0.000}{1.959}\right)^2 + \left(\frac{0.00025}{0.01925}\right)^2} = 0.0005\text{ N}\cdot\text{m} \\
\alpha &= \frac{0.00954}{0.01925} = 0.496\text{ rad/s}^2 \\
\delta\alpha &= 0.496\sqrt{\left(\frac{0.00086}{0.00954}\right)^2 + \left(\frac{0.00025}{0.01925}\right)^2} = 0.045\text{ rad/s}^2 \\
I &= \frac{0.0377}{0.496} = 0.0761\text{ kg}\cdot\text{m}^2 \\
\delta I &= 0.0761\sqrt{\left(\frac{0.0005}{0.0377}\right)^2 + \left(\frac{0.045}{0.496}\right)^2} = 0.0070\text{ kg}\cdot\text{m}^2 \\
I_{\text{masses}} &= 0.0761 - 0.0122 = 0.0639\text{ kg}\cdot\text{m}^2 \\
\delta I_{\text{masses}} &= \sqrt{0.0070^2 + 0.0027^2} = 0.0075\text{ kg}\cdot\text{m}^2
\end{align*}

\textbf{Trial 20 — $R = 0.302 \pm 0.001\text{ m}$:}
\begin{align*}
a &= \frac{2(0.590)}{(12.24)^2} = 0.00788\text{ m/s}^2 \\
\delta a &= 0.00788\sqrt{\left(\frac{0.001}{0.590}\right)^2 + \left(\frac{2(0.5)}{12.24}\right)^2} = 0.00064\text{ m/s}^2 \\
g - a &= 9.807 - 0.00788 = 9.799 \pm 0.001\text{ m/s}^2 \\
T &= 0.2(9.799) = 1.960\text{ N}, \qquad \delta T = 0.2(0.00064) = 0.000\text{ N} \\
\tau &= (1.960)(0.01925) = 0.0377\text{ N}\cdot\text{m} \\
\delta\tau &= 0.0377\sqrt{\left(\frac{0.000}{1.960}\right)^2 + \left(\frac{0.00025}{0.01925}\right)^2} = 0.0005\text{ N}\cdot\text{m} \\
\alpha &= \frac{0.00788}{0.01925} = 0.409\text{ rad/s}^2 \\
\delta\alpha &= 0.409\sqrt{\left(\frac{0.00064}{0.00788}\right)^2 + \left(\frac{0.00025}{0.01925}\right)^2} = 0.034\text{ rad/s}^2 \\
I &= \frac{0.0377}{0.409} = 0.0922\text{ kg}\cdot\text{m}^2 \\
\delta I &= 0.0922\sqrt{\left(\frac{0.0005}{0.0377}\right)^2 + \left(\frac{0.034}{0.409}\right)^2} = 0.0077\text{ kg}\cdot\text{m}^2 \\
I_{\text{masses}} &= 0.0922 - 0.0122 = 0.0800\text{ kg}\cdot\text{m}^2 \\
\delta I_{\text{masses}} &= \sqrt{0.0077^2 + 0.0027^2} = 0.0082\text{ kg}\cdot\text{m}^2
\end{align*}

\textbf{Trial 21 — $R = 0.312 \pm 0.001\text{ m}$:}
\begin{align*}
a &= \frac{2(0.590)}{(13.64)^2} = 0.00634\text{ m/s}^2 \\
\delta a &= 0.00634\sqrt{\left(\frac{0.001}{0.590}\right)^2 + \left(\frac{2(0.5)}{13.64}\right)^2} = 0.00047\text{ m/s}^2 \\
g - a &= 9.807 - 0.00634 = 9.801 \pm 0.001\text{ m/s}^2 \\
T &= 0.2(9.801) = 1.960\text{ N}, \qquad \delta T = 0.2(0.00047) = 0.000\text{ N} \\
\tau &= (1.960)(0.01925) = 0.0377\text{ N}\cdot\text{m} \\
\delta\tau &= 0.0377\sqrt{\left(\frac{0.000}{1.960}\right)^2 + \left(\frac{0.00025}{0.01925}\right)^2} = 0.0005\text{ N}\cdot\text{m} \\
\alpha &= \frac{0.00634}{0.01925} = 0.329\text{ rad/s}^2 \\
\delta\alpha &= 0.329\sqrt{\left(\frac{0.00047}{0.00634}\right)^2 + \left(\frac{0.00025}{0.01925}\right)^2} = 0.025\text{ rad/s}^2 \\
I &= \frac{0.0377}{0.329} = 0.1145\text{ kg}\cdot\text{m}^2 \\
\delta I &= 0.1145\sqrt{\left(\frac{0.0005}{0.0377}\right)^2 + \left(\frac{0.025}{0.329}\right)^2} = 0.0087\text{ kg}\cdot\text{m}^2 \\
I_{\text{masses}} &= 0.1145 - 0.0122 = 0.1023\text{ kg}\cdot\text{m}^2 \\
\delta I_{\text{masses}} &= \sqrt{0.0087^2 + 0.0027^2} = 0.0091\text{ kg}\cdot\text{m}^2
\end{align*}

\begin{table}[h!]
\centering
\small
\begin{tabular}{cccc}
\toprule
\textbf{Trial} & \textbf{$R$ (m)} & \textbf{$I_{\text{total}}$ ($\text{kg}\cdot\text{m}^2$)} & \textbf{$I_{\text{masses}}$ ($\text{kg}\cdot\text{m}^2$)} \\
\midrule
2 & $0.122 \pm 0.001$ & $0.018 \pm 0.003$ & $0.006 \pm 0.004$ \\
3 & $0.132 \pm 0.001$ & $0.021 \pm 0.004$ & $0.009 \pm 0.005$ \\
4 & $0.142 \pm 0.001$ & $0.024 \pm 0.004$ & $0.012 \pm 0.005$ \\
5 & $0.152 \pm 0.001$ & $0.028 \pm 0.004$ & $0.016 \pm 0.005$ \\
6 & $0.162 \pm 0.001$ & $0.027 \pm 0.004$ & $0.015 \pm 0.005$ \\
7 & $0.172 \pm 0.001$ & $0.030 \pm 0.004$ & $0.018 \pm 0.005$ \\
8 & $0.182 \pm 0.001$ & $0.035 \pm 0.005$ & $0.023 \pm 0.006$ \\
9 & $0.192 \pm 0.001$ & $0.038 \pm 0.005$ & $0.026 \pm 0.006$ \\
10 & $0.202 \pm 0.001$ & $0.041 \pm 0.005$ & $0.029 \pm 0.006$ \\
11 & $0.212 \pm 0.001$ & $0.043 \pm 0.005$ & $0.031 \pm 0.006$ \\
12 & $0.222 \pm 0.001$ & $0.048 \pm 0.006$ & $0.036 \pm 0.007$ \\
13 & $0.232 \pm 0.001$ & $0.050 \pm 0.006$ & $0.038 \pm 0.007$ \\
14 & $0.242 \pm 0.001$ & $0.054 \pm 0.006$ & $0.042 \pm 0.007$ \\
15 & $0.252 \pm 0.001$ & $0.055 \pm 0.006$ & $0.043 \pm 0.007$ \\
16 & $0.262 \pm 0.001$ & $0.060 \pm 0.006$ & $0.048 \pm 0.007$ \\
17 & $0.272 \pm 0.001$ & $0.0665 \pm 0.0065$ & $0.0543 \pm 0.0070$ \\
18 & $0.282 \pm 0.001$ & $0.0712 \pm 0.0067$ & $0.0590 \pm 0.0072$ \\
19 & $0.292 \pm 0.001$ & $0.0761 \pm 0.0070$ & $0.0639 \pm 0.0075$ \\
20 & $0.302 \pm 0.001$ & $0.0922 \pm 0.0077$ & $0.0800 \pm 0.0082$ \\
21 & $0.312 \pm 0.001$ & $0.1145 \pm 0.0087$ & $0.1023 \pm 0.0091$ \\
\bottomrule
\end{tabular}
\end{table}

\subsection*{Linearization}
The graph of $I_{\text{masses}}$ vs.\ $R$ is not linear. It is modeled as a power function $I = kR^j$, which is linearized by taking the natural log of both sides:
\[
\ln(I) = j\ln(R) + \ln(k)
\]
A graph of $\ln(I)$ vs.\ $\ln(R)$ is linear with slope $j$ and y-intercept $\ln(k)$.
\[
\delta(\ln R) = \frac{\delta R}{R}, \qquad \delta(\ln I) = \frac{\delta I}{I}
\]

\begin{flushleft}
\small
$\ln(0.122) = -2.104, \quad \delta = \frac{0.001}{0.122} = 0.008 \qquad\qquad \ln(0.006) = -5.12, \quad \delta = \frac{0.004}{0.006} = 0.67$ \\[4pt]
$\ln(0.132) = -2.025, \quad \delta = \frac{0.001}{0.132} = 0.008 \qquad\qquad \ln(0.009) = -4.71, \quad \delta = \frac{0.005}{0.009} = 0.56$ \\[4pt]
$\ln(0.142) = -1.952, \quad \delta = \frac{0.001}{0.142} = 0.007 \qquad\qquad \ln(0.012) = -4.42, \quad \delta = \frac{0.005}{0.012} = 0.42$ \\[4pt]
$\ln(0.152) = -1.884, \quad \delta = \frac{0.001}{0.152} = 0.007 \qquad\qquad \ln(0.016) = -4.14, \quad \delta = \frac{0.005}{0.016} = 0.31$ \\[4pt]
$\ln(0.162) = -1.820, \quad \delta = \frac{0.001}{0.162} = 0.006 \qquad\qquad \ln(0.015) = -4.20, \quad \delta = \frac{0.005}{0.015} = 0.33$ \\[4pt]
$\ln(0.172) = -1.760, \quad \delta = \frac{0.001}{0.172} = 0.006 \qquad\qquad \ln(0.018) = -4.02, \quad \delta = \frac{0.005}{0.018} = 0.28$ \\[4pt]
$\ln(0.182) = -1.704, \quad \delta = \frac{0.001}{0.182} = 0.005 \qquad\qquad \ln(0.023) = -3.77, \quad \delta = \frac{0.006}{0.023} = 0.26$ \\[4pt]
$\ln(0.192) = -1.650, \quad \delta = \frac{0.001}{0.192} = 0.005 \qquad\qquad \ln(0.026) = -3.65, \quad \delta = \frac{0.006}{0.026} = 0.23$ \\[4pt]
$\ln(0.202) = -1.599, \quad \delta = \frac{0.001}{0.202} = 0.005 \qquad\qquad \ln(0.029) = -3.54, \quad \delta = \frac{0.006}{0.029} = 0.21$ \\[4pt]
$\ln(0.212) = -1.551, \quad \delta = \frac{0.001}{0.212} = 0.005 \qquad\qquad \ln(0.031) = -3.47, \quad \delta = \frac{0.006}{0.031} = 0.19$ \\[4pt]
$\ln(0.222) = -1.505, \quad \delta = \frac{0.001}{0.222} = 0.005 \qquad\qquad \ln(0.036) = -3.32, \quad \delta = \frac{0.007}{0.036} = 0.19$ \\[4pt]
$\ln(0.232) = -1.461, \quad \delta = \frac{0.001}{0.232} = 0.004 \qquad\qquad \ln(0.038) = -3.27, \quad \delta = \frac{0.007}{0.038} = 0.18$ \\[4pt]
$\ln(0.242) = -1.419, \quad \delta = \frac{0.001}{0.242} = 0.004 \qquad\qquad \ln(0.042) = -3.17, \quad \delta = \frac{0.007}{0.042} = 0.17$ \\[4pt]
$\ln(0.252) = -1.378, \quad \delta = \frac{0.001}{0.252} = 0.004 \qquad\qquad \ln(0.043) = -3.15, \quad \delta = \frac{0.007}{0.043} = 0.16$ \\[4pt]
$\ln(0.262) = -1.339, \quad \delta = \frac{0.001}{0.262} = 0.004 \qquad\qquad \ln(0.048) = -3.04, \quad \delta = \frac{0.007}{0.048} = 0.15$ \\[4pt]
$\ln(0.272) = -1.302, \quad \delta = \frac{0.001}{0.272} = 0.004 \qquad\qquad \ln(0.0543) = -2.91, \quad \delta = \frac{0.0070}{0.0543} = 0.13$ \\[4pt]
$\ln(0.282) = -1.266, \quad \delta = \frac{0.001}{0.282} = 0.004 \qquad\qquad \ln(0.0590) = -2.83, \quad \delta = \frac{0.0072}{0.0590} = 0.12$ \\[4pt]
$\ln(0.292) = -1.231, \quad \delta = \frac{0.001}{0.292} = 0.003 \qquad\qquad \ln(0.0639) = -2.75, \quad \delta = \frac{0.0075}{0.0639} = 0.12$ \\[4pt]
$\ln(0.302) = -1.197, \quad \delta = \frac{0.001}{0.302} = 0.003 \qquad\qquad \ln(0.0800) = -2.53, \quad \delta = \frac{0.0082}{0.0800} = 0.10$ \\[4pt]
$\ln(0.312) = -1.165, \quad \delta = \frac{0.001}{0.312} = 0.003 \qquad\qquad \ln(0.1023) = -2.28, \quad \delta = \frac{0.0091}{0.1023} = 0.09$
\end{flushleft}

\begin{table}[h!]
\centering
\small
\caption{Linearized Data}
\begin{tabular}{ccc}
\toprule
\textbf{Trial} & \textbf{$\ln(R)$} & \textbf{$\ln(I)$} \\
\midrule
2 & $-2.104 \pm 0.008$ & $-5.12 \pm 0.67$ \\
3 & $-2.025 \pm 0.008$ & $-4.71 \pm 0.56$ \\
4 & $-1.952 \pm 0.007$ & $-4.42 \pm 0.42$ \\
5 & $-1.884 \pm 0.007$ & $-4.14 \pm 0.31$ \\
6 & $-1.820 \pm 0.006$ & $-4.20 \pm 0.33$ \\
7 & $-1.760 \pm 0.006$ & $-4.02 \pm 0.28$ \\
8 & $-1.704 \pm 0.005$ & $-3.77 \pm 0.26$ \\
9 & $-1.650 \pm 0.005$ & $-3.65 \pm 0.23$ \\
10 & $-1.599 \pm 0.005$ & $-3.54 \pm 0.21$ \\
11 & $-1.551 \pm 0.005$ & $-3.47 \pm 0.19$ \\
12 & $-1.505 \pm 0.005$ & $-3.32 \pm 0.19$ \\
13 & $-1.461 \pm 0.004$ & $-3.27 \pm 0.18$ \\
14 & $-1.419 \pm 0.004$ & $-3.17 \pm 0.17$ \\
15 & $-1.378 \pm 0.004$ & $-3.15 \pm 0.16$ \\
16 & $-1.339 \pm 0.004$ & $-3.04 \pm 0.15$ \\
17 & $-1.302 \pm 0.004$ & $-2.91 \pm 0.13$ \\
18 & $-1.266 \pm 0.004$ & $-2.83 \pm 0.12$ \\
19 & $-1.231 \pm 0.003$ & $-2.75 \pm 0.12$ \\
20 & $-1.197 \pm 0.003$ & $-2.53 \pm 0.10$ \\
21 & $-1.165 \pm 0.003$ & $-2.28 \pm 0.09$ \\
\bottomrule
\end{tabular}
\end{table}

% Uncomment and insert your graph filenames when ready:
% \begin{figure}[h!]
%     \centering
%     \includegraphics[width=0.75\textwidth]{raw_graph.png}
%     \caption{Moment of Inertia of Adjustable Masses ($I_{\text{masses}}$) vs.\ Radial Distance ($R$).}
% \end{figure}
%
% \begin{figure}[h!]
%     \centering
%     \includegraphics[width=0.75\textwidth]{linearized_graph.png}
%     \caption{Linearized Graph: $\ln(I)$ vs.\ $\ln(R)$.}
% \end{figure}

\textbf{Best-Fit Line (through $(-2.025, -4.71)$ and $(-1.165, -2.28)$):}
\begin{align*}
\Delta\ln(I) &= -2.28 - (-4.71) = 2.43 \\
\Delta\ln(R) &= -1.165 - (-2.025) = 0.860 \\
j &= \frac{2.43}{0.860} = 2.83 \\
\ln(k) &= -4.71 - (2.83)(-2.025) = -4.71 + 5.72 = 1.01 \\
\ln(I) &= 2.83\ln(R) + 1.01
\end{align*}

\textbf{Maximum Slope Line (through $(-2.025, -5.27)$ and $(-1.165, -2.19)$):}
\begin{align*}
-4.71 - 0.56 &= -5.27, \qquad -2.28 + 0.09 = -2.19 \\
j_{\max} &= \frac{-2.19 - (-5.27)}{0.860} = \frac{3.08}{0.860} = 3.58 \\
\ln(k)_{\max} &= -5.27 - (3.58)(-2.025) = -5.27 + 7.25 = 1.98
\end{align*}

\textbf{Minimum Slope Line (through $(-2.025, -4.15)$ and $(-1.165, -2.37)$):}
\begin{align*}
-4.71 + 0.56 &= -4.15, \qquad -2.28 - 0.09 = -2.37 \\
j_{\min} &= \frac{-2.37 - (-4.15)}{0.860} = \frac{1.78}{0.860} = 2.07 \\
\ln(k)_{\min} &= -4.15 - (2.07)(-2.025) = -4.15 + 4.19 = 0.04
\end{align*}

\textbf{Uncertainties \& Final Relationship:}
\begin{align*}
\delta j &= \frac{j_{\max} - j_{\min}}{2} = \frac{3.58 - 2.07}{2} = 0.76 \implies j = 2.8 \pm 0.8 \\
\delta\ln(k) &= \frac{|\ln(k)_{\max} - \ln(k)_{\min}|}{2} = \frac{|1.98 - 0.04|}{2} = 0.97 \implies \ln(k) = 1.0 \pm 1.0 \\
k &= e^{1.01} = 2.75\text{ kg}\cdot\text{m}^{-0.8}, \qquad \delta k = k\,\delta\ln(k) = 2.75(0.97) = 2.67 \implies k = 2.8 \pm 2.7\text{ kg}\cdot\text{m}^{-0.8} \\
I &= (2.8 \pm 2.7)R^{2.8 \pm 0.8}
\end{align*}

\section*{Conclusion}
Moving the adjustable masses farther from the axis of rotation increased their moment of inertia, confirming that rotational inertia depends on how mass is distributed relative to the axis rather than total mass alone. The experimental relationship found was $I = (2.8 \pm 2.7)R^{2.8 \pm 0.8}$, where $I$ is in $\text{kg}\cdot\text{m}^2$ and $R$ is in $\text{m}$. For point masses, $I = (4m)R^2$ predicts an exponent of $2$. The measured exponent range of $2.0$ to $3.6$ barely reaches $2$ at its lower bound, with a central value steeper than expected. Systematic and human error explain this discrepancy: kinetic friction in the demonstrator's bearings exerts a retarding torque that increases drop times $t$, and because $I \propto t^2$ with the longest drops occurring at larger radii, friction disproportionately inflates $I$ at outer positions and steepens the slope. Additionally, treating the adjustable cylinders as point masses ignores their own rotational inertia about their centers of mass ($I = 4(I_{\text{cm}} + mR^2)$), while reaction time during hand-timing introduces random scatter across trials.

\end{document}